\documentclass[11pt,a4paper]{article}

\usepackage[margin=2.5cm]{geometry}
\usepackage[T1]{fontenc}
\usepackage{lmodern}
\usepackage{orcidlink}

\usepackage[backend=biber, style=numeric, sorting=none]{biblatex}

\addbibresource{ref.bib}

\title{\textbf{Closing the Gap Between Proof and Implementation:\\
Formalizing Elligator in Lean}}

\author{
Chris Anto Fröschl\orcidlink{0009-0006-3562-5395}\\
Munich University of Applied Sciences\\
Department of Computer Science and Mathematics\\
Email: chris.froeschl@hm.edu
}

\date{}

\begin{document}

\maketitle

\section{Introduction}

Elliptic Curve Cryptography is widely used in modern secure systems, yet its security
guarantees typically rely on paper proofs that are disconnected from concrete implementations.
This separation makes it difficult to ensure that theoretical results are faithfully realized
in practice.

We present a formalization of Elligator \cite{bernstein2013a}
in the Lean theorem prover \cite{demoura2021a}. Elligator encodes elliptic curve points as bit strings
indistinguishable from uniform random data, enabling censorship-resistant key exchange.
Our development provides a machine-checked implementation together with formal proofs
of its correctness and indistinguishability properties.

\section{Formalization Approach}

The formalization covers the full stack of required mathematics, including finite field
arithmetic, elliptic curve operations, and the encoding and decoding procedures underlying
Elligator. Translating the original proof into a formal setting requires making implicit
assumptions explicit, restructuring arguments into smaller verifiable steps, and resolving
edge cases that are often omitted in informal presentations.

In doing so, the original construction is not only reproduced but made fully explicit,
yielding a precise and unambiguous specification of both the algorithm and its guarantees.
This process highlights subtle dependencies between algebraic properties and correctness
arguments that are not immediately visible in the paper proof.

\section{Lean and Mathematical Foundations}

The work builds on Lean’s mathematical library Mathlib \cite{mathlib2019a}, leveraging existing formalizations
of algebraic structures and elliptic curves to ensure a robust and reusable foundation.
It also aligns with the emerging CSLib ecosystem~\cite{demoura2026a}, which aims to provide
a shared library of formalized computer science theory for the verification of software
systems and cryptographic protocols.

\section{Contribution and Impact}

To the best of our knowledge, this constitutes the first fully machine-checked implementation
of Elligator.

This result demonstrates how formal methods can directly connect cryptographic theory with
verified implementations of software and protocols. Moreover, it showcases that modern theorem
provers are capable of handling formalizations of contemporary research mathematics at the level
of current cryptographic constructions, strengthening confidence in their applicability to
security-critical systems.

\printbibliography

\end{document}
